\chapter{Instrumento y estructura de la encuesta}
\label{apx:encuesta}

Las preguntas y opciones se reproducen conforme a la copia documental del formulario aplicado. Se conservaron las formulaciones originales del instrumento, incluidas sus inconsistencias de tratamiento (por ejemplo, ``experimentaste'' y ``quieras'' frente al ``usted'' del resto). El diccionario de variables y la tabla de escalas presentan las denominaciones normalizadas utilizadas para la organización y el análisis de los datos.

\section*{1. Alcance y estructura del instrumento}
\addcontentsline{toc}{subsection}{1. Alcance y estructura del instrumento}
El formulario original se encontraba ramificado por rol e incluía recorridos diferenciados para estudiantes, docentes, gestión académica y perfiles técnicos. En este apéndice se reproduce únicamente el recorrido correspondiente a estudiantes y la sección final común. El análisis cuantitativo del \textit{User Research} utilizó exclusivamente la rama estudiantil.

\begin{center}\small
\begin{tabular}{@{}p{5.2cm}p{3.2cm}p{5.5cm}@{}}
\toprule
\rowcolor{accent!15}\textbf{Componente} & \textbf{Cantidad} & \textbf{Carácter documentado} \\
\midrule
Pregunta inicial de segmentación por rol & 1 pregunta & Obligatoria \\
Rama estudiantes & 17 preguntas & 16 obligatorias y 1 opcional \\
Sección final común & 5 preguntas & 4 obligatorias y 1 opcional \\
\bottomrule
\end{tabular}
\end{center}
\subsection*{PARTE I --- REPRODUCCIÓN DEL INSTRUMENTO APLICADO}
\section*{2. Pregunta inicial de segmentación por rol}
\addcontentsline{toc}{subsection}{2. Pregunta inicial de segmentación por rol}
La respuesta correspondiente a ``Estudiante universitario/a o terciario/a'' derivaba al recorrido estudiantil reproducido a continuación.

\begin{infobox}\small
\textbf{\textcolor{accent}{Código:}} ROL\\
\textbf{\textcolor{accent}{Pregunta:}} ¿Desde qué rol responde esta encuesta?\\
\textbf{\textcolor{accent}{Tipo de pregunta:}} Opción única\\
\textbf{\textcolor{accent}{Carácter:}} Obligatoria\\
\textbf{\textcolor{accent}{Opciones de respuesta:}} • Estudiante universitario/a o terciario/a \\ • Docente, ayudante o evaluador/a \\ • Coordinador/a, directivo/a o responsable académico \\ • Perfil técnico, campus virtual, sistemas o tecnología educativa \\ • Other: [opción de respuesta libre]\\
\end{infobox}
\section*{3. Rama estudiantes --- EST01 a EST17}
\addcontentsline{toc}{subsection}{3. Rama estudiantes --- EST01 a EST17}
La rama estudiantil comprende 17 preguntas: 16 obligatorias y una pregunta abierta opcional (EST17).

\begin{infobox}\small
\textbf{\textcolor{accent}{Código:}} EST01\\
\textbf{\textcolor{accent}{Pregunta:}} ¿En qué área o carrera estudia o estudió?\\
\textbf{\textcolor{accent}{Tipo de pregunta:}} Opción única\\
\textbf{\textcolor{accent}{Carácter:}} Obligatoria\\
\textbf{\textcolor{accent}{Opciones de respuesta:}} • Ingeniería / Tecnología / Informática \\ • Ciencias Económicas / Administración / Negocios \\ • Derecho / Ciencias Sociales \\ • Comunicación / Diseño \\ • Salud \\ • Educación \\ • Other: [opción de respuesta libre]\\
\end{infobox}
\begin{infobox}\small
\textbf{\textcolor{accent}{Código:}} EST02\\
\textbf{\textcolor{accent}{Pregunta:}} ¿En qué etapa aproximada de la carrera se encuentra?\\
\textbf{\textcolor{accent}{Tipo de pregunta:}} Opción única\\
\textbf{\textcolor{accent}{Carácter:}} Obligatoria\\
\textbf{\textcolor{accent}{Opciones de respuesta:}} • Primer año \\ • Segundo año \\ • Tercer año \\ • Cuarto año \\ • Quinto año o más \\ • Ya finalicé la carrera \\ • No aplica\\
\end{infobox}
\begin{infobox}\small
\textbf{\textcolor{accent}{Código:}} EST03\\
\textbf{\textcolor{accent}{Pregunta:}} ¿Rindió alguna evaluación digital o remota durante los últimos dos años?\\
\textbf{\textcolor{accent}{Tipo de pregunta:}} Opción única\\
\textbf{\textcolor{accent}{Carácter:}} Obligatoria\\
\textbf{\textcolor{accent}{Opciones de respuesta:}} • Sí \\ • No \\ • No recuerdo\\
\end{infobox}
\begin{infobox}\small
\textbf{\textcolor{accent}{Código:}} EST04\\
\textbf{\textcolor{accent}{Pregunta:}} ¿Qué tipos de evaluaciones digitales o remotas realizó?\\
\textbf{\textcolor{accent}{Tipo de pregunta:}} Selección múltiple\\
\textbf{\textcolor{accent}{Carácter:}} Obligatoria\\
\textbf{\textcolor{accent}{Opciones de respuesta:}} • Cuestionarios \textit{online} con opciones múltiples \\ • Parciales escritos \textit{online} \\ • Finales \textit{online} \\ • Evaluaciones orales por videollamada \\ • Entregas asincrónicas de trabajos \\ • Evaluaciones en campus virtual \\ • No realicé evaluaciones digitales o remotas \\ • Other: [opción de respuesta libre]\\
\end{infobox}
\begin{infobox}\small
\textbf{\textcolor{accent}{Código:}} EST05\\
\textbf{\textcolor{accent}{Pregunta:}} En general, ¿qué tan confiables considera que son las evaluaciones digitales para reflejar el conocimiento real del estudiante?\\
\textbf{\textcolor{accent}{Tipo de pregunta:}} Escala lineal\\
\textbf{\textcolor{accent}{Carácter:}} Obligatoria\\
\textbf{\textcolor{accent}{Opciones de respuesta:}} • Nada confiables \\ • 1 \\ • 2 \\ • 3 \\ • 4 \\ • 5 \\ • Muy confiables\\
\end{infobox}
\begin{infobox}\small
\textbf{\textcolor{accent}{Código:}} EST06\\
\textbf{\textcolor{accent}{Pregunta:}} En comparación con una evaluación presencial, ¿considera que en una evaluación digital es más difícil garantizar condiciones equitativas para todos los estudiantes?\\
\textbf{\textcolor{accent}{Tipo de pregunta:}} Opción única\\
\textbf{\textcolor{accent}{Carácter:}} Obligatoria\\
\textbf{\textcolor{accent}{Opciones de respuesta:}} • Sí \\ • No \\ • Depende del tipo de evaluación \\ • No estoy seguro/a\\
\end{infobox}
\begin{infobox}\small
\textbf{\textcolor{accent}{Código:}} EST07\\
\textbf{\textcolor{accent}{Pregunta:}} ¿Qué situaciones considera que pueden afectar la integridad académica en una evaluación digital?\\
\textbf{\textcolor{accent}{Tipo de pregunta:}} Selección múltiple\\
\textbf{\textcolor{accent}{Carácter:}} Obligatoria\\
\textbf{\textcolor{accent}{Opciones de respuesta:}} • Consulta de materiales no permitidos \\ • Ayuda de otra persona durante la evaluación \\ • Comunicación entre estudiantes \\ • Suplantación de identidad \\ • Búsqueda de respuestas en internet \\ • Problemas técnicos que afectan el control de la evaluación \\ • Falta de claridad sobre las reglas \\ • No considero que haya situaciones relevantes \\ • Other: [opción de respuesta libre]\\
\end{infobox}
\begin{infobox}\small
\textbf{\textcolor{accent}{Código:}} EST08\\
\textbf{\textcolor{accent}{Pregunta:}} ¿Qué mecanismos de supervisión o control experimentaste o conoce en evaluaciones digitales?\\
\textbf{\textcolor{accent}{Tipo de pregunta:}} Selección múltiple\\
\textbf{\textcolor{accent}{Carácter:}} Obligatoria\\
\textbf{\textcolor{accent}{Opciones de respuesta:}} • Cámara encendida durante la evaluación \\ • Micrófono encendido \\ • Pantalla compartida \\ • Grabación de la sesión \\ • Bloqueo de navegador \\ • Control de tiempo \\ • Preguntas aleatorias o mezcladas \\ • Revisión posterior del docente \\ • Registro de actividades en campus virtual \\ • No experimenté mecanismos de supervisión \\ • Other: [opción de respuesta libre]\\
\end{infobox}
\begin{infobox}\small
\textbf{\textcolor{accent}{Código:}} EST09\\
\textbf{\textcolor{accent}{Pregunta:}} ¿Qué tan invasivo le parece el uso obligatorio de cámara durante toda una evaluación digital?\\
\textbf{\textcolor{accent}{Tipo de pregunta:}} Escala lineal\\
\textbf{\textcolor{accent}{Carácter:}} Obligatoria\\
\textbf{\textcolor{accent}{Opciones de respuesta:}} • Nada invasivo \\ • 1 \\ • 2 \\ • 3 \\ • 4 \\ • 5 \\ • Muy invasivo\\
\end{infobox}
\begin{infobox}\small
\textbf{\textcolor{accent}{Código:}} EST10\\
\textbf{\textcolor{accent}{Pregunta:}} ¿Qué tan invasivo le parece compartir pantalla durante toda una evaluación digital?\\
\textbf{\textcolor{accent}{Tipo de pregunta:}} Escala lineal\\
\textbf{\textcolor{accent}{Carácter:}} Obligatoria\\
\textbf{\textcolor{accent}{Opciones de respuesta:}} • Nada invasivo \\ • 1 \\ • 2 \\ • 3 \\ • 4 \\ • 5 \\ • Muy invasivo\\
\end{infobox}
\begin{infobox}\small
\textbf{\textcolor{accent}{Código:}} EST11\\
\textbf{\textcolor{accent}{Pregunta:}} ¿Qué tan invasivo le parecería que una plataforma registre eventos de interacción, como uso de teclado, mouse o cambios de pestaña, sin registrar el contenido escrito por el estudiante?\\
\textbf{\textcolor{accent}{Tipo de pregunta:}} Escala lineal\\
\textbf{\textcolor{accent}{Carácter:}} Obligatoria\\
\textbf{\textcolor{accent}{Opciones de respuesta:}} • Nada invasivo \\ • 1 \\ • 2 \\ • 3 \\ • 4 \\ • 5 \\ • Muy invasivo\\
\end{infobox}
\begin{infobox}\small
\textbf{\textcolor{accent}{Código:}} EST12\\
\textbf{\textcolor{accent}{Pregunta:}} ¿Qué información consideraría razonable que una institución registre durante una evaluación digital?\\
\textbf{\textcolor{accent}{Tipo de pregunta:}} Selección múltiple\\
\textbf{\textcolor{accent}{Carácter:}} Obligatoria\\
\textbf{\textcolor{accent}{Opciones de respuesta:}} • Hora de inicio y finalización \\ • Tiempo dedicado a cada pregunta o sección \\ • Cambios de pestaña o ventana \\ • Eventos de teclado y mouse sin registrar el contenido escrito \\ • Dirección IP o dispositivo utilizado \\ • Grabación de cámara \\ • Grabación de pantalla \\ • Contenido exacto escrito por el estudiante \\ • Ninguna de las anteriores \\ • Other: [opción de respuesta libre]\\
\end{infobox}
\begin{infobox}\small
\textbf{\textcolor{accent}{Código:}} EST13\\
\textbf{\textcolor{accent}{Pregunta:}} ¿Qué información consideraría excesiva o invasiva durante una evaluación digital?\\
\textbf{\textcolor{accent}{Tipo de pregunta:}} Selección múltiple\\
\textbf{\textcolor{accent}{Carácter:}} Obligatoria\\
\textbf{\textcolor{accent}{Opciones de respuesta:}} • Cámara encendida durante toda la evaluación \\ • Micrófono encendido durante toda la evaluación \\ • Grabación de pantalla completa \\ • Registro del contenido escrito \\ • Registro de eventos de teclado y mouse \\ • Registro de cambios de pestaña o ventana \\ • Dirección IP o dispositivo utilizado \\ • Ninguna de las anteriores \\ • Other: [opción de respuesta libre]\\
\end{infobox}
\begin{infobox}\small
\textbf{\textcolor{accent}{Código:}} EST14\\
\textbf{\textcolor{accent}{Pregunta:}} Si un sistema detectara una situación atípica durante una evaluación, ¿qué debería ocurrir?\\
\textbf{\textcolor{accent}{Tipo de pregunta:}} Opción única\\
\textbf{\textcolor{accent}{Carácter:}} Obligatoria\\
\textbf{\textcolor{accent}{Opciones de respuesta:}} • Debería generarse una alerta para revisión humana posterior \\ • Debería notificarse automáticamente al estudiante \\ • Debería anularse automáticamente la evaluación \\ • Debería solicitarse una instancia adicional de validación \\ • No debería utilizarse ningún sistema automático \\ • No estoy seguro/a\\
\end{infobox}
\begin{infobox}\small
\textbf{\textcolor{accent}{Código:}} EST15\\
\textbf{\textcolor{accent}{Pregunta:}} ¿Qué tan importante considera que una alerta generada por un sistema NO implique automáticamente una acusación contra el estudiante?\\
\textbf{\textcolor{accent}{Tipo de pregunta:}} Escala lineal\\
\textbf{\textcolor{accent}{Carácter:}} Obligatoria\\
\textbf{\textcolor{accent}{Opciones de respuesta:}} • Nada importante \\ • 1 \\ • 2 \\ • 3 \\ • 4 \\ • 5 \\ • Muy importante\\
\end{infobox}
\begin{infobox}\small
\textbf{\textcolor{accent}{Código:}} EST16\\
\textbf{\textcolor{accent}{Pregunta:}} ¿Qué condiciones harían más aceptable el uso de herramientas de apoyo a la supervisión en evaluaciones digitales?\\
\textbf{\textcolor{accent}{Tipo de pregunta:}} Selección múltiple\\
\textbf{\textcolor{accent}{Carácter:}} Obligatoria\\
\textbf{\textcolor{accent}{Opciones de respuesta:}} • Que se informe previamente qué datos se registran \\ • Que exista consentimiento del estudiante \\ • Que no se almacene contenido sensible \\ • Que las alertas sean revisadas por una persona \\ • Que el sistema no tome decisiones automáticas de sanción \\ • Que se expliquen los motivos de una alerta \\ • Que los datos se eliminen después de cierto tiempo \\ • Que exista un canal de reclamo o revisión \\ • Ninguna condición lo haría aceptable \\ • Other: [opción de respuesta libre]\\
\end{infobox}
\begin{infobox}\small
\textbf{\textcolor{accent}{Código:}} EST17\\
\textbf{\textcolor{accent}{Pregunta:}} Desde su experiencia, ¿cuál es el principal problema de las evaluaciones digitales o remotas?\\
\textbf{\textcolor{accent}{Tipo de pregunta:}} Respuesta abierta\\
\textbf{\textcolor{accent}{Carácter:}} Opcional\\
\textbf{\textcolor{accent}{Opciones de respuesta:}} ---\\
\end{infobox}
\section*{4. Sección final común --- COM01 a COM05}
\addcontentsline{toc}{subsection}{4. Sección final común --- COM01 a COM05}
La sección final común comprende 5 preguntas: 4 obligatorias y una pregunta abierta opcional (COM05).

\begin{infobox}\small
\textbf{\textcolor{accent}{Código:}} COM01\\
\textbf{\textcolor{accent}{Pregunta:}} ¿Qué tan importante es que la persona evaluada conozca previamente qué datos se registran y con qué finalidad?\\
\textbf{\textcolor{accent}{Tipo de pregunta:}} Escala lineal\\
\textbf{\textcolor{accent}{Carácter:}} Obligatoria\\
\textbf{\textcolor{accent}{Opciones de respuesta:}} • Nada importante \\ • 1 \\ • 2 \\ • 3 \\ • 4 \\ • 5 \\ • Muy importante\\
\end{infobox}
\begin{infobox}\small
\textbf{\textcolor{accent}{Código:}} COM02\\
\textbf{\textcolor{accent}{Pregunta:}} ¿Qué tan importante es que exista revisión humana antes de tomar una decisión académica basada en una alerta?\\
\textbf{\textcolor{accent}{Tipo de pregunta:}} Escala lineal\\
\textbf{\textcolor{accent}{Carácter:}} Obligatoria\\
\textbf{\textcolor{accent}{Opciones de respuesta:}} • Nada importante \\ • 1 \\ • 2 \\ • 3 \\ • 4 \\ • 5 \\ • Muy importante\\
\end{infobox}
\begin{infobox}\small
\textbf{\textcolor{accent}{Código:}} COM03\\
\textbf{\textcolor{accent}{Pregunta:}} ¿Qué tan aceptable le parece que una herramienta registre eventos de interacción, como teclado, mouse o navegación, si no almacena el contenido escrito?\\
\textbf{\textcolor{accent}{Tipo de pregunta:}} Escala lineal\\
\textbf{\textcolor{accent}{Carácter:}} Obligatoria\\
\textbf{\textcolor{accent}{Opciones de respuesta:}} • Nada aceptable \\ • 1 \\ • 2 \\ • 3 \\ • 4 \\ • 5 \\ • Muy aceptable\\
\end{infobox}
\begin{infobox}\small
\textbf{\textcolor{accent}{Código:}} COM04\\
\textbf{\textcolor{accent}{Pregunta:}} ¿Qué condición considera más importante para que una herramienta de apoyo a la supervisión sea justa?\\
\textbf{\textcolor{accent}{Tipo de pregunta:}} Opción única\\
\textbf{\textcolor{accent}{Carácter:}} Obligatoria\\
\textbf{\textcolor{accent}{Opciones de respuesta:}} • Transparencia sobre los datos recolectados \\ • Consentimiento informado \\ • Revisión humana \\ • No almacenar contenido sensible \\ • Explicabilidad de las alertas \\ • Posibilidad de reclamo o revisión \\ • Baja tasa de falsos positivos \\ • Other: [opción de respuesta libre]\\
\end{infobox}
\begin{infobox}\small
\textbf{\textcolor{accent}{Código:}} COM05\\
\textbf{\textcolor{accent}{Pregunta:}} ¿Hay algún comentario adicional que quieras agregar sobre evaluaciones digitales, privacidad o integridad académica?\\
\textbf{\textcolor{accent}{Tipo de pregunta:}} Respuesta abierta\\
\textbf{\textcolor{accent}{Carácter:}} Opcional\\
\textbf{\textcolor{accent}{Opciones de respuesta:}} ---\\
\end{infobox}
\subsection*{PARTE II --- DOCUMENTACIÓN NORMALIZADA PARA EL ANÁLISIS}
Las tablas de esta parte reproducen la documentación normalizada de la planilla de trabajo (\textit{workbook}) definitiva utilizada para organizar las variables y las escalas del análisis. Los códigos de variable (por ejemplo, \texttt{EST03\_evaluacion\_digital\_ultimos\_2\_anios}) se escriben sin tildes ni eñes porque son los nombres de columna del conjunto de datos y de los \textit{scripts} de análisis; la pregunta original, con su ortografía completa, figura en la columna contigua.

\section*{5. Diccionario de variables}
\addcontentsline{toc}{subsection}{5. Diccionario de variables}
\begin{center}\footnotesize
\begin{longtable}{@{}p{3.4cm}p{3.7cm}p{3.7cm}p{3.7cm}@{}}
\toprule
\rowcolor{accent!15}\textbf{Código} & \textbf{Pregunta original} & \textbf{Tipo} & \textbf{Origen} \\
\midrule\endhead
EST01\_\allowbreak area\_\allowbreak carrera & ¿En qué área o carrera estudia o estudió? & Opción única & Rama estudiantes \\
EST02\_\allowbreak etapa\_\allowbreak carrera & ¿En qué etapa aproximada de la carrera se encuentra? & Opción única & Rama estudiantes \\
EST03\_\allowbreak evaluacion\_\allowbreak digital\_\allowbreak ultimos\_\allowbreak 2\_\allowbreak anios & ¿Rindió alguna evaluación digital o remota durante los últimos dos años? & Opción única & Rama estudiantes \\
EST04\_\allowbreak tipos\_\allowbreak evaluaciones & ¿Qué tipos de evaluaciones digitales o remotas realizó? & Selección múltiple & Rama estudiantes \\
EST05\_\allowbreak confiabilidad\_\allowbreak evaluacion\_\allowbreak digital & En general, ¿qué tan confiables considera que son las evaluaciones digitales para reflejar el conocimiento real del estudiante? & Escala lineal & Rama estudiantes \\
EST06\_\allowbreak equidad\_\allowbreak vs\_\allowbreak presencial & En comparación con una evaluación presencial, ¿considera que en una evaluación digital es más difícil garantizar condiciones equitativas para todos los estudiantes? & Opción única & Rama estudiantes \\
EST07\_\allowbreak situaciones\_\allowbreak integridad & ¿Qué situaciones considera que pueden afectar la integridad académica en una evaluación digital? & Selección múltiple & Rama estudiantes \\
EST08\_\allowbreak mecanismos\_\allowbreak supervision & ¿Qué mecanismos de supervisión o control experimentó o conoce en evaluaciones digitales? & Selección múltiple & Rama estudiantes \\
EST09\_\allowbreak invasividad\_\allowbreak camara & ¿Qué tan invasivo le parece el uso obligatorio de cámara durante toda una evaluación digital? & Escala lineal & Rama estudiantes \\
EST10\_\allowbreak invasividad\_\allowbreak compartir\_\allowbreak pantalla & ¿Qué tan invasivo le parece compartir pantalla durante toda una evaluación digital? & Escala lineal & Rama estudiantes \\
EST11\_\allowbreak invasividad\_\allowbreak eventos\_\allowbreak interaccion & ¿Qué tan invasivo le parecería que una plataforma registre eventos de interacción, como uso de teclado, mouse o cambios de pestaña, sin registrar el contenido escrito por el estudiante? & Escala lineal & Rama estudiantes \\
EST12\_\allowbreak informacion\_\allowbreak razonable & ¿Qué información consideraría razonable que una institución registre durante una evaluación digital? & Selección múltiple & Rama estudiantes \\
EST13\_\allowbreak informacion\_\allowbreak excesiva & ¿Qué información consideraría excesiva o invasiva durante una evaluación digital? & Selección múltiple & Rama estudiantes \\
EST14\_\allowbreak accion\_\allowbreak ante\_\allowbreak situacion\_\allowbreak atipica & Si un sistema detectara una situación atípica durante una evaluación, ¿qué debería ocurrir? & Opción única & Rama estudiantes \\
EST15\_\allowbreak importancia\_\allowbreak no\_\allowbreak acusacion\_\allowbreak automatica & ¿Qué tan importante considera que una alerta generada por un sistema NO implique automáticamente una acusación contra el estudiante? & Escala lineal & Rama estudiantes \\
EST16\_\allowbreak condiciones\_\allowbreak aceptabilidad & ¿Qué condiciones harían más aceptable el uso de herramientas de apoyo a la supervisión en evaluaciones digitales? & Selección múltiple & Rama estudiantes \\
EST17\_\allowbreak principal\_\allowbreak problema & Desde su experiencia, ¿cuál es el principal problema de las evaluaciones digitales o remotas? & Respuesta abierta & Rama estudiantes \\
COM01\_\allowbreak transparencia\_\allowbreak datos & ¿Qué tan importante es que la persona evaluada conozca previamente qué datos se registran y con qué finalidad? & Escala lineal & Sección final común \\
COM02\_\allowbreak revision\_\allowbreak humana & ¿Qué tan importante es que exista revisión humana antes de tomar una decisión académica basada en una alerta? & Escala lineal & Sección final común \\
COM03\_\allowbreak aceptabilidad\_\allowbreak registro\_\allowbreak eventos & ¿Qué tan aceptable le parece que una herramienta registre eventos de interacción, como teclado, mouse o navegación, si no almacena el contenido escrito? & Escala lineal & Sección final común \\
COM04\_\allowbreak condicion\_\allowbreak justicia & ¿Qué condición considera más importante para que una herramienta de apoyo a la supervisión sea justa? & Opción única & Sección final común \\
COM05\_\allowbreak comentario\_\allowbreak adicional & ¿Hay algún comentario adicional que quiera agregar sobre evaluaciones digitales, privacidad o integridad académica? & Respuesta abierta & Sección final común \\
\bottomrule
\end{longtable}
\end{center}
\section*{6. Escalas 1--5 y anclajes verificados}
\addcontentsline{toc}{subsection}{6. Escalas 1--5 y anclajes verificados}
Los anclajes se presentan con las denominaciones normalizadas verificadas en la hoja ``Escalas''. Las escalas miden constructos distintos ---confiabilidad, invasividad, importancia y aceptabilidad--- y no se combinan en un índice único.

\begin{center}\footnotesize
\begin{longtable}{@{}p{3.2cm}p{2.1cm}p{1.6cm}p{2.1cm}p{2.1cm}p{2.6cm}@{}}
\toprule
\rowcolor{accent!15}\textbf{Código} & \textbf{Constructo} & \textbf{Escala} & \textbf{Ancla 1} & \textbf{Ancla 5} & \textbf{Interpretación} \\
\midrule\endhead
EST05\_\allowbreak confiabilidad\_\allowbreak evaluacion\_\allowbreak digital & Confiabilidad & 1--5 & Nada confiable & Muy confiable & Un valor más alto indica mayor confiabilidad percibida. \\
EST09\_\allowbreak invasividad\_\allowbreak camara & Invasividad & 1--5 & Nada invasivo & Muy invasivo & Un valor más alto indica mayor invasividad percibida. \\
EST10\_\allowbreak invasividad\_\allowbreak compartir\_\allowbreak pantalla & Invasividad & 1--5 & Nada invasivo & Muy invasivo & Un valor más alto indica mayor invasividad percibida. \\
EST11\_\allowbreak invasividad\_\allowbreak eventos\_\allowbreak interaccion & Invasividad & 1--5 & Nada invasivo & Muy invasivo & Un valor más alto indica mayor invasividad percibida. \\
EST15\_\allowbreak importancia\_\allowbreak no\_\allowbreak acusacion\_\allowbreak automatica & Importancia & 1--5 & Nada importante & Muy importante & Un valor más alto indica mayor importancia asignada. \\
COM01\_\allowbreak transparencia\_\allowbreak datos & Importancia & 1--5 & Nada importante & Muy importante & Un valor más alto indica mayor importancia asignada. \\
COM02\_\allowbreak revision\_\allowbreak humana & Importancia & 1--5 & Nada importante & Muy importante & Un valor más alto indica mayor importancia asignada. \\
COM03\_\allowbreak aceptabilidad\_\allowbreak registro\_\allowbreak eventos & Aceptabilidad & 1--5 & Nada aceptable & Muy aceptable & Un valor más alto indica mayor aceptabilidad percibida. \\
\bottomrule
\end{longtable}
\end{center}